\scanpage{69}

% top, connector
donc

% top, Hom and defining condition
\[
\operatorname{Hom}(A^n_*,C^{n\,!}_{*M})=
\{A^n_*\xrightarrow{f^{\unsure{*|n}}}C^n_{*M}\mid \sigma^*\circ f=0\quad\forall\sigma:\Delta_n\to\Delta_m\text{ avec }m<n\}
\]

% upper left, small square
\margin{\diagram{\resizebox{\linewidth}{!}{\begin{tikzcd}[ampersand replacement=\&,column sep=small,row sep=small]
A^n_\bullet\arrow[r,"f^n"]\arrow[d,"\sigma_A^*"'] \& C^n_*\arrow[d,"\sigma_C^*"]\\
A^m_\bullet\arrow[r] \& C^m_*
\end{tikzcd}}}}

% upper, given map
Or \deleted{\unsure{A}} $A^n_\bullet\mathrel{\unsure{\relbar|\longrightarrow}}C^n_\bullet$ est donné par

% upper, map labelled K
\[
A^n_n\xrightarrow{\unsure{K_n|k_n|\kappa_n}}M
\]

% upper, composite
son composé $A^n_\bullet\relbar C^n_\bullet\relbar C^m_*$ est donné par \deleted{la}

% upper middle, M-valued map
\[
A^n_m\longrightarrow M
\]

% upper middle, description
qui n'est autre que le composé

% upper middle, factorisation
\[
A^n_m\xrightarrow{\sigma^*}A^n_n\xrightarrow{K_n}M
\]

% middle, cancellation and new opening
\deleted{et la condition $f(A^n$\unsure{)}} on a donc

% middle, criterion and bound
\[
f^n(A^n_*)\subset C^n_\bullet\quad\Longleftrightarrow\quad
\left|\begin{aligned}
&K_n\circ\sigma^*=0\\
&\forall\sigma:\Delta_n\longrightarrow\Delta_m,\\[-2pt]
&\hspace{2em}\underline{m<n}
\end{aligned}\right.
\]

% middle, degeneracies and reference to a)
(il suffit $=$ de \unsure{l'écrire|la vérifier|le vérifier} quand $\sigma^*$ est une \unsure{op.|opér.} de dégénérescence --) ; or cette condition est tjrs satisfaite si $A_m^n=0$ si $m<n$, ce qu'on a supposé dans a).

% lower, example
\uline{Ex} Soit $S$ un anneau à \unsure{pd-id.|idéal|p.d.} $J$ à p.d., $t\in J$, considérons alors \unsure{l'an$^{\mathrm t}$|cet|l'} Anneau différentiel $\frac{1}{2}$ simplicial à p.d.

% lower, divided-power quotient
\[
A^*_{*S}=(\Gamma^*(\Phi_{*S})\otimes_S\Lambda^*\Psi_{*S})/(T-t.)_{\mathrm{pd}},
\]

% lower, resolution and integrals
c'est une \uline{résolution} de $S$, et $A_j^i=0$ si $i>j$, donc on a un syst. d'intégrales de Stokes-De Rham à valeurs dans $S$, d'où

% bottom, isomorphism and diagonal arrow
\diagram{\begin{tikzcd}[ampersand replacement=\&,column sep=small,row sep=small]
A^*_\bullet\arrow[r,"\sim"] \& \bigl(\Gamma^*(\Phi_{*S})\otimes_{S}\Lambda^*\Psi_{*S}/(T-t.)_{\mathrm{pd}}\bigr)\arrow[dr]\\
\& \& \Lambda^{*+1}_{S}\Phi_{*S}
\end{tikzcd}}

% bottom edge, final remark
dont il faudrait étudier les propriétés multiplicatives\deleted{\unsure{app.}}\ldots

% lower left
% archive sheet number: 68
